% ============================================================
% ch30.tex —— 第 30 章 交点算术：从 2×2=4 到贝祖定理
% ============================================================
\chapter{交点算术：从 \texorpdfstring{$2\times2=4$}{2×2=4} 到贝祖定理}

上一章末尾挂出的猜想：$m$ 次曲线与 $n$ 次曲线的交点恰好 $mn$ 个（复数世界、计重数）。这一章把它当实验做——数据会先支持、再出卖、最后成全这个猜想。出卖之处，正是代数几何埋宝藏的地方。

\section{实验：交点乘法表}

\paragraph{实验一：两条抛物线，真能交四点吗？}直接试 $y = x^2$ 与 $y = -x^2 + 2x + 2$（一条开口向上、一条向下）：联立 $x^2 = -x^2 + 2x + 2$，即 $x^2 - x - 1 = 0$——两个 $x$，各给一个交点。只有 $2$ 个！四点在哪？

别急着改猜想——\textbf{这两条抛物线一上一下，图形只在中间地带重叠}，四点凑不齐情有可原。要凑四个交点，得找"你穿过来、我穿过去"的姿态。经典构造：\textbf{横卧抛物线 $\times$ 竖立抛物线}——$y^2 = x$ 与 $y = x^2$。代入 $y = x^2$：
\[
x^4 = x \;\Longrightarrow\; x^4 - x = x(x^3 - 1) = x(x-1)(x^2 + x + 1) = 0 .
\]
四个根：$x = 0$、$x = 1$、以及 $x^2+x+1 = 0$ 的两个复根 $\omega = \tfrac{-1+\sqrt3\mathrm{i}}2$ 与 $\omega^2 = \tfrac{-1-\sqrt3\mathrm{i}}2$（三次单位根，即 $1$ 的三个立方根中除 $1$ 外的两个）。对应四个交点：
\[
(0,0),\quad (1,1),\quad \left(\omega, \omega^2\right),\quad \left(\omega^2, \omega\right).
\]
（验算 $(1,1)$：两条曲线都过它 \checkmark；$(0,0)$ 同 \checkmark。）前两个看得见（实数），后两个住在复数空间——\textbf{四票到齐}，$2\times2 = 4$。猜想的第一张选票，含两票"看不见的选票"。

\paragraph{实验二：切点怎么算？}$y = x^2$ 与 $y = 0$（横轴）：相切于原点，看图数点得 $1$，而 $1\times2 = 2$——\textbf{猜想阵亡？}先别下葬。看代数：联立即 $x^2 = 0$，\textbf{二重根}。若切点算一票，交点数 $1$，破账；若\textbf{切点算两票}，$2 = 1\times2$，账平。这不是和稀泥的补救——判别式为零时根的重数本来就是二（第 29 章），是"\textbf{看图数点}"数错了，代数从来没骗人。\textbf{计重数规则}：图像只是代数的取景框，代数说两票就是两票。（第 29 章选读小节"动根滑回定根"的实验已经演示了这个"重合"的过程：割线的两个交点随直线转动逐渐合拢，重合为切点——票没有消失，是叠在了一起。）

\paragraph{实验三：平行线叛逃。}$y = 0$ 与 $y = 1$：两条平行直线，永不相交——实数交点 $0$，复数交点也 $0$（两方程无公共解，$y$ 不可能既是 $0$ 又是 $1$，与 $x$ 虚实无关）。而 $1\times1 = 1$。\textbf{这次连复数都救不了：猜想真的破了一个洞。}

三幕剧的前两幕演完：数据支持（实验一）、数据出卖（实验二被计重数救回，实验三彻底破洞）。破洞的地方藏着金矿。

\section{补洞：把"无穷远"设为合法住址}

洞怎么补？老思路（第 29 章方法论）：\textbf{交点不会凭空消失，只会搬家}。平行线的交点搬去哪了？——请出画家。

文艺复兴的画家面对的技术难题：怎么把三维世界"忠实"画到二维画布上？他们的答案（透视法）里有一条你天天见的现象：\textbf{平行的铁轨、走廊边线、林荫道两侧，在画面里汇聚于同一个"消失点"}。平行（现实）与相交（画面）并不矛盾——隔着无穷远，它们相交于画面的消失点。画家的约定俗成，被几何学家扶正：

\begin{itemize}
\item 给每一个\textbf{方向}配一个\textbf{无穷远点}（水平方向一个、竖直方向一个、东偏北 $17^\circ$ 方向一个……同向的平行线共享同一个）；
\item 规定：平行线恰在它们方向的无穷远点\textbf{相交}（记一票——两条不同直线的交点本来就是一票，账目恰好）；
\item 全体无穷远点手拉手组成一条\textbf{无穷远直线}（它是所有方向的"花名册"）。
\end{itemize}

\begin{tikzpicture}[scale=0.9]
  \draw[->] (-4.6,0) -- (4.6,0);
  \draw[thick,shenlan] (-4.2,-1.5) -- (4.2,1.5);
  \draw[thick,molv] (-4.2,-2.6) -- (4.2,0.4);
  \draw[dashed,cheng] (4.35,1.55) circle (0.13);
  \node[cheng!70!black] at (4.55,2.25) {\scriptsize 同一方向};
  \node[cheng!70!black] at (4.55,1.8) {\scriptsize 的无穷远点};
  \node[shenlan] at (-3.5,1.35) {\scriptsize 铁轨 $1$};
  \node[molv] at (-3.5,-3.05) {\scriptsize 铁轨 $2$};
  \node at (0,-3.65) {\scriptsize 平行线在"消失点"相遇——画家用了五百年的知识};
\end{tikzpicture}

普通平面 $\cup$ 无穷远直线 $=$ 升级版空间。在升级版里，\textbf{任何两条不同直线恰好交于一点}：相交的交在老地方，平行的交在无穷远。$y=0$ 与 $y=1$ 的那一票，记在"水平方向的无穷远点"名下——实验三的洞补上了。

\begin{faming}
\textbf{土名}：画家平面 $\to$ \textbf{正式名}：\textbf{射影平面}（projective plane）。它是普通平面加上所有方向的无穷远点及一条无穷远直线。"射影"取自投影（projection）：画家把三维世界投影到画布上，平行线的像相交——画布本身就是一张天然的射影平面。
\end{faming}

\section{定理宣读}

补好两个洞（无穷远点、计重数），猜想从"实验规律"升格为定理：

\begin{dingli}[贝祖定理，1779]
在复数射影平面上，不共享整条曲线的 $m$ 次曲线与 $n$ 次曲线，\textbf{恰好}交于 $mn$ 个点（交点计重数，无穷远点也算）。
\end{dingli}

全部账目对清（含两处前面欠的悬案）：

\begin{center}
\begin{tabular}{lcc}
\toprule
案例 & 看图数点 & 贝祖记账 \\
\midrule
直线 $\times$ 圆（一般位置） & $2$ & $1\times2 = 2$ \\
直线 $\times$ 圆（相切） & $1$ & 切点两票：$2 = 1\times2$ \\
直线 $\times$ 圆（实数离弦） & $0$ & 复数两点：$2 = 1\times2$ \\
两平行直线 & $0$ & 无穷远点一票：$1 = 1\times1$ \\
横轴 $\times$ 双曲线 $xy=1$ & $0$ & 无穷远\textbf{重数二}一票：$2 = 1\times2$ \\
两条抛物线（横$\times$竖） & $2$（实） & 复数四点：$4 = 2\times2$ \\
\bottomrule
\end{tabular}
\end{center}

倒数第二行是第 29 章留下的悬案，也是"计重数"最精彩的演出：横轴与双曲线 $xy = 1$ 在水平无穷远点相遇时，\textbf{贴得比相切还紧}（重数二）——渐近线那个"无限逼近"的执念（第 29 章）在无穷远处兑现为"二阶接触"。这笔账的亲手验算放在第 31 章选读小节：五分钟齐次坐标，一算便知。

\section{次数：曲线世界的座次表}

贝祖定理让"次数"从记号变成了权力的度量：\textbf{次数就是交点额度}。三次曲线与直线必有 $3\times1 = 3$ 个交点（计重、含复、含无穷远）——第 32 章的椭圆曲线（三次）将把这个 $3$ 当信用卡额度刷出惊世骇俗的操作：过曲线上两点作直线，必交\textbf{第三点}（两票已占，第三票\textbf{被迫存在}）。\textbf{曲线上长出群运算}的整场好戏，幕布就是贝祖定理拉开的。

\begin{lishi}
贝祖（\'Etienne B\'ezout，1730--1783，法国）1779 年的著作《代数方程通论》里给出的是交点数的\textbf{上界}版本（且在实数范围内陈述）；"复数射影平面上恰好 $mn$ 个（计重数）"的现代版，是十九世纪射影几何与复几何成熟后回填的，教科书惯例整体归功于他，本书从俗。"重数"的严格代数定义（交局部环的维数）要等到二十世纪；本书用的"切点两票、渐近三阶接触"是几何直觉版，足够看戏。
\end{lishi}

\begin{dongshou}
\begin{itemize}
  \yisi 数交点（计重数、复数）：$y = x^2$ 与 $x = y^2$（第 29 章实验一，四点 $2$ 实 $2$ 复）；$y = x^3$ 与 $y = x$（$x^3 = x$，$x(x-1)(x+1) = 0$，三点 $3\times1 = 3$ 全实——看图：三次曲线与直线三次穿叉）。
  \yier 构造两条二次曲线恰好交于 $4$ 个实点：圆 $x^2 + y^2 = 2$ 与双曲线 $xy = \tfrac12$。代入 $y = \tfrac{1}{2x}$：$x^2 + \tfrac1{4x^2} = 2$，即 $4x^4 - 8x^2 + 1 = 0$，$x^2 = 1 \pm\tfrac{\sqrt3}{2}$，两个正解各开正负平方——\textbf{四个实交点}，全部可见。画草图核对。
  \yier 用贝祖额度秒杀一道作图题："过椭圆外一点能作椭圆的几条切线？"（切线 $=$ 直线与二次曲线"交点重合"的情形，两票叠一处。过外点的直线束里，与椭圆重票相交的恰有 $2$ 条——请用第 29 章选读的"动根追踪法"验证：直线 $y - y_0 = m(x - x_0)$ 代入 $x^2 + 4y^2 = 4$，得关于 $x$ 的二次方程，令 $\Delta(m) = 0$——$\Delta$ 是 $m$ 的二次式，恰两根，即两条切线。）
\end{itemize}
\end{dongshou}

\begin{shaonao}
\textbf{贝祖定理的"逆"：用自由度数点。}平面上给定 $5$ 个一般位置的点（无三点共线），恰有\textbf{一条}二次曲线穿过它们。账这么算：一般二次曲线 $Ax^2 + Bxy + Cy^2 + Dx + Ey + F = 0$ 有 $6$ 个系数，但整体乘一个非零常数是同一条曲线——\textbf{自由度 $5$}；每个点提供一条方程（代入坐标），$5$ 点 $5$ 方程，解唯一确定（一般位置时）。照方抓药：

(1) 过 $9$ 个一般位置的点，有几条三次曲线？（自由度 $10 - 1 = 9$，一般解唯一——\textbf{恰好一条}。这是第 32 章"弦切游戏"的地基：三次曲线被 $9$ 个点几乎钉死，所以"两点连线交出第三点"的构造才这么刚性。）

(2) 过 $4$ 个点呢？（自由度 $5 > 4$，欠定——\textbf{一整族}（单参数族）二次曲线都能过这 $4$ 点。想象 $4$ 颗钉子，能挂住一整排形状各异的橡皮筋。）

(3) 帕斯卡定理预告：圆内接六边形三组对边交点\textbf{共线}——用"两条三次曲线过同样九点"加贝祖定理可以三行证出（牛顿的思路）。这是"自由度 $+$ 贝祖"组合拳的名场面，完整版留给大学，先记下这道门牌。
\end{shaonao}

\begin{chengshi}
"$\omega$、$\omega^2$ 是三次单位根"用到了复数开方（$x^3 = 1$ 除 $x=1$ 外的两个根），中学没学但代入验证即可：$\omega^3 = \omega\cdot\omega^2 = \omega\cdot\tfrac{-1-\sqrt3\mathrm{i}}2$……直接乘开约两行，请自行验算。贝祖定理的完整证明需要结式（消元的代数机器，判断两个多项式有无公共根的行列式判据）与射影坐标——放行大学；"重数"的直觉定义（切点两票、渐近二阶接触）与严格定义（局部交环维数）在光滑情形下吻合。两条曲线"共享整条分支"（如 $xy = 0$ 与 $x = 0$）时定理明文除外——那时交点无穷多，账目爆表。
\end{chengshi}

洞已补齐，定理已立。下一章：走进"无穷远点"这个新大陆本身——先拿画家五百年前的直觉当门票，再用五分钟的齐次坐标，亲手把第 29、30 章的每笔"悬案"算成铁案。
